\section{詳細版について}

本書は、技術報告書の詳細版である。技術報告書の内容に加えて、次のものを収めた。

\begin{itemize}
\item \textbf{基礎編}（第 I 部の前）：剛体力学の基礎（第~\ref{ch:basics-rigid}章）、幾何力学入門（第~\ref{ch:basics-geom}章）、
  最適化・制御・学習の基礎（第~\ref{ch:basics-opt}章）。
\item \textbf{各章の詳細解説}：各章の末尾に、式の詳しい導出、数値を使った例題、アルゴリズムの手順、結果の全データ、
  追加の解析（感度分析など）、考察課題を置いた。
\item \textbf{記号表}（目次の直後）と\textbf{用語集}（付録）。
\end{itemize}

\subsection{読み方の案内}

\begin{longtable}{p{0.30\linewidth}p{0.62\linewidth}}
\toprule
読者 & 推奨する読み方 \\ \midrule \endhead
結論を知りたい & 序章 → 各章の本文（詳細解説は飛ばす）→ 結論 \\
力学の背景から理解したい & 序章 → 基礎編 → 第 I 部から順に \\
特定のテーマに関心がある & 序章 → 関心のある部（第 I〜III 部）→ その章の詳細解説 \\
再現・拡張したい & 付録（数値計算の方法、パラメータ、再現手順）と各章の詳細解説 \\
\bottomrule
\end{longtable}

\subsection{研究の系譜}

手足の動きで体の向きを変える問題は、AMBAC の登場（1981 年 \cite{R42}）より前から研究されてきた。
落下する猫が足から着地する現象は、19 世紀末の写真による観察以来の古い問題で、1969 年に Kane と Scher が 2 つの円柱をつないだモデルで力学的に説明した \cite{R7}。
翌 1970 年には、同じ著者らが宇宙飛行士が手足の動きで体の向きを変える方法を解析している \cite{R8}。
1980〜90 年代には、Montgomery らがこれらを幾何学（ゲージ理論）の言葉で統一的に理解し \cite{R4}、
同じ時期に宇宙ロボットの分野で、アームの動きが母機の姿勢を乱す問題と、それを利用・抑制する制御が研究された \cite{R12,R13,R14,R15}。
本報告書は、この流れにフィクションの AMBAC を重ね、さらに地上の安定性・人の肩・学習へとつなげる試みである。
